\section{本章小结}
\label{sec:cprob-summary}

本章用一个盒子贯穿始终：$3$个黑球、$2$个白球，摸出两个。把这次抽样从头算一遍，就能把本章的概念串起来，并看清各步依赖的假设。表~\ref{tab:cprob-summary}按顺序列出各步。

\begin{BookInlineTable}
  \centering
  \caption{摸球问题中各步的计算与所用假设}
  \label{tab:cprob-summary}
  \begin{tabular}{lll}
    \toprule
    步骤 & 计算 & 所用假设 \\
    \midrule
    样本空间 & 有序对共$20$个 & 记录先后，细到能表达事件 \\
    等可能 & 每个有序对概率$1/20$ & 搅匀后每个球机会相同 \\
    计数 & $\mathbb{P}(H_1\cap H_2)=6/20$ & 分子分母同一口径 \\
    条件概率 & $\mathbb{P}(H_2\mid H_1)=2/4$ & $\mathbb{P}(H_1)>0$ \\
    乘法规则 & $\tfrac35\cdot\tfrac24=\tfrac{3}{10}$ & 条件随已摸结果改变 \\
    全概率 & $\mathbb{P}(H_2)=\tfrac{6}{20}+\tfrac{6}{20}=\tfrac35$ & $H_1,H_1^{\mathrm{c}}$构成划分 \\
    Bayes & $\mathbb{P}(H_1\mid H_2)=1/2$ & $\mathbb{P}(H_2)>0$ \\
    独立性 & 不放回不独立，放回独立 & 放回时乘积空间等可能 \\
    \bottomrule
  \end{tabular}
\end{BookInlineTable}

\paragraph{样本空间与计数}
概率首先要有一个样本空间：它列出一次试验的全部可能结果，事件是其中的子集，一次试验只产生一个结果，却可以让许多事件同时发生。样本空间由记录方式决定，它必须细到能表达所关心的事件。在有限且等可能的前提下，事件的概率是计数比$\lvert A\rvert/\lvert\Omega\rvert$，它满足非负、归一和互斥可加三条性质，补事件公式、加法公式与联合界都从这三条推出。

等可能是对试验的假设，依据是物理的对称性，不是样本点名称的对称。点数和、颜色记录、无序的可重复组合都合并了不同数量的基本结果，不能直接数类别。计数时分子和分母必须在同一个样本空间中数；放回与不放回改变的是样本空间本身。相对频率可以检验等可能假设是否与试验相符，但本章不用它定义概率。

\paragraph{条件与独立}
条件概率把样本范围缩小到已知发生的事件$B$，在其中按原比例重新归一，要求$\mathbb{P}(B)>0$。由它得到乘法规则，可以沿试验过程逐步计算；按划分分组得到全概率公式，再反推来源得到Bayes公式。先验、条件检出概率和后验分别回答不同的问题，$\mathbb{P}(A\mid B)$与$\mathbb{P}(B\mid A)$不能互换。逐次更新时，前一次的后验成为下一次的先验，但每次乘以似然比都要求给定状态时各次观察互不影响。

独立是交事件概率等于乘积这一数量关系，它与互斥不同，也不由“物理上是否有关”决定。多个事件的相互独立要求所有子组的乘积关系，两两独立或只验证全体之交都不够。重复试验中的独立性来自乘积样本空间上的等可能假设，不放回抽样不满足这一假设。条件之下的独立与不加条件的独立是两回事。

\paragraph{边界与下一步}
古典模型严格而具体，但只适用于有限且等可能的结果。结果不等可能时要逐个指定概率；结果有可数无穷多个时，可加性必须推广到可数个事件；结果取连续位置时，单点的概率为零，区间的概率要由长度或密度给出。下一章把试验的结果记录成数值，引入随机变量，用概率表、分布函数和密度描述它取各种值的概率，这三类情形都在那里得到具体的处理。
